

\ch{Appendix D: Institutional Review Board Human Subjects Application}\label{appendix:irb}

This appendix reproduces the Institutional Review Board (IRB) Human Subjects Application submitted to the Dakota State University IRB for this study. The application records the protocol as reviewed: its purpose, participant tasks, recruitment method, target population, privacy and data-retention safeguards, and consent mechanism. Consistent with the Ethical Considerations discussion in Chapter~3, the study collected no personally identifiable information and received an exempt review determination (DSUIRB-20260421-13EX, April 21, 2026) before any recruitment or data collection began. The responses below are reproduced as submitted. The principal investigator's personal telephone number has been withheld for privacy.

\section{Investigators and Administrative Details}

\begin{table}[H]
  \centering
  \caption{IRB application administrative details.}\label{tab:irb-admin}
  % ALT: Administrative header of the IRB application: project title, submission date 27 March 2026, principal investigator Patrick Engebretson and co-investigator John Bonar (both Beacom College of Computer and Cyber Sciences), co-investigator CITI training 22 January 2025, and co-investigator signature dated 20 April 2026.
  \begin{tabular}{@{}p{4cm}p{9.2cm}@{}}
    \toprule
    \textbf{Field} & \textbf{Value} \\
    \midrule
    Project title & Transforming Information Assurance and IT Service Management Through Digital Engineering \\
    Date submitted & 27 March 2026 \\
    Principal Investigator & Patrick Engebretson, Beacom College of Computer and Cyber Sciences (\texttt{pat.engebretson@dsu.edu}) \\
    Co-Investigator & John Bonar, Beacom College of Computer and Cyber Sciences (\texttt{john.bonar@trojans.dsu.edu}) \\
    CITI training (Co-Investigator) & 22 January 2025 \\
    Co-Investigator signature date & 20 April 2026 \\
    Faculty supervisor & Patrick Engebretson \\
    \bottomrule
  \end{tabular}
\end{table}

\section{Application Responses}

\begin{enumerate}[leftmargin=2.4em,itemsep=2pt]
  \item \textbf{Potential conflict of interest.} No.
  \item \textbf{External funding.} No.
  \item \textbf{Required to earn a DSU academic degree.} Yes. The dissertation and its Dissertation Proposal Defense Approval were provided as the required supporting documentation.
  \item \textbf{Submitted to another IRB.} No.
  \item \textbf{Non-IRB approval letters (committees, schools, or external organizations).} Not applicable.
  \item \textbf{Purpose of the research.} Identify perceptions of applicability and value among Information Technology and Information Assurance professionals with regard to Digital Engineering.
  \item \textbf{Tasks participants perform.} Complete a six-part survey containing 31 questions in total.\footnote{One item was subsequently identified as a duplicate and removed before administration, so the instrument as administered contained thirty items (Appendix~A).}
  \item \textbf{Contacting participants and protecting privacy.} Recruit through LinkedIn posts to a professional social-media network and through email solicitation to industry peers: (1) a LinkedIn post, reproduced in Appendix~A, and (2) an email containing the solicitation text. No identifying information is collected, which preserves participant privacy.
  \item \textbf{Maximum participants enrolled.} A minimum of 385 completed responses are required to achieve a margin of error of 5\%.
  \item \textbf{Vulnerable study populations.} None. The study targeted no children, prisoners, pregnant persons, non-English speakers, students, employees of the investigators, or other listed vulnerable groups.
  \item \textbf{Targeting a specific ethnic group.} Not applicable.
  \item \textbf{Age range of subjects.} 18 to 64 years.
  \item \textbf{Inclusion and exclusion criteria.} Inclusion is specific to invited Information Assurance and Information Technology professionals.
  \item \textbf{Research location.} Online participation within the United States.
  \item \textbf{Participant compensation.} No.
  \item \textbf{Direct benefits to participants.} No direct benefit.
  \item \textbf{Private information collected.} None. The study collected no names, contact information, device or serial identifiers, geographic detail, dates, records, or other private information.
  \item \textbf{Traceability of a record to a participant.} No.
  \item \textbf{Agreement to destroy identifying data at study conclusion.} Yes.
  \item \textbf{Agreement to retain de-identified data for three years.} Yes.
  \item \textbf{Associated risks (physical, economic, psychological, legal, social, or use of private records).} None of the listed risks apply.
  \item \textbf{Consent type.} Consent Statement for adults, permissible because no identifiers are collected.
  \item \textbf{Method of obtaining consent.} Web-based.
  \item \textbf{HIPAA authorization or waiver.} Not applicable.
\end{enumerate}

The co-investigator signed the application on 20 April 2026. Faculty supervisor Dr.~Patrick Engebretson provided the supervisor assurance that the research is scientifically and scholarly sound, that competencies and resources are adequate, and that no conflicts of interest exist.
